\section{From pinned returns to spontaneous magnetization}
\label{sec:pressure}

The return estimates hold for every deterministic choice of extra pins.
A positive magnetic field is an exact mixture of such choices. We use
this fact to bound finite-volume pressure increments, then take the
thermodynamic limit before differentiating at zero field.

\subsection{Pressure and the cost of boundary pins}

Write $V=|\Lambda_N|$ and
\[
 p_N(h)=\frac1{\beta V}\log Z_{\Lambda_N}(h).
\]
The limit $p_\beta(h)=\lim_{N\to\infty}p_N(h)$ exists for every
$\beta>0$ and $h\in\R$. Here is the finite-range argument. Tile a large
cube, apart from a remainder, by cubes of a fixed side $l$ and delete
all interactions between tiles. There are $O(V/l)$ deleted bonds,
each with uniformly bounded operator norm; the remainder has $o(V)$
sites as $N\to\infty$ with $l$ fixed and may be left as singletons.
If two Hamiltonians differ by an operator of norm at most $a$, their
logarithmic partition functions differ by at most $\beta a$, by
ordered-eigenvalue comparison. Consequently the limiting upper and
lower pressures both lie within $C/l$ of the pressure per site of a
single tile. Letting $l\to\infty$ proves existence. Since
$[H^J_{S,N},M_{S,N}]=0$ and $\norm{M_{S,N}}\le SV$, each $p_N$ is
convex and $S$-Lipschitz in $h$. These properties pass to $p_\beta$,
so its right derivative at zero exists.

Let $Z^+_\Lambda(h)$ be the loop expression in
Proposition~\ref{prop:loop-representation} with all cycles meeting
$D_0$ constrained to have color $+1$, and put
$p_N^+(h)=(\beta V)^{-1}\log Z^+_{\Lambda_N}(h)$.

\begin{proposition}[Boundary pressure equivalence]
\label{prop:pressure-equivalence}
For every fixed $\beta>0$ and $h\ge0$,
\begin{equation}
 \lim_{N\to\infty}p_N^+(h)=p_\beta(h).
 \label{eq:boundary-pressure}
\end{equation}
\end{proposition}

\begin{proof}
Factor $e^{\beta hSV}$ from the unrestricted color sum. Its picture
weight, relative to the base law, becomes
\[
 W_h=\prod_C(1+e^{-\beta h a_C}).
\]
Given the picture, the colors of different cycles are independent and
the probability that a cycle is up is at least $1/2$. If $b_\partial$
counts the cycles meeting $D_0$, the corresponding normalized picture
expectation therefore satisfies
\begin{equation}
 1\ge \frac{Z^+_\Lambda(h)}{Z_\Lambda(h)}
 \ge \mathbb E_h[2^{-b_\partial}]
 \ge \exp\{-\log(2)\,\mathbb E_h b_\partial\}.
 \label{eq:surface-cost}
\end{equation}
A cycle visiting the boundary either contains a time-zero boundary
slot or arrives there at a mark incident to the boundary. Thus
\[
 b_\partial\le \ell|\partial\Lambda|+2N_\partial,
\]
where $N_\partial$ is the number of such marks. Insertion of a mark
changes at most two cycle factors in $W_h$, each between one and two,
so it increases $W_h$ by at most a factor of four. The Poisson insertion
argument of Lemma~\ref{lem:factorial} gives
$\mathbb E_hN_\partial\le C\beta|\partial\Lambda|$.
Taking logarithms in \eqref{eq:surface-cost} yields
\[
 0\le p_N(h)-p_N^+(h)
 \le C(1+\beta)\frac{|\partial\Lambda|}{\beta|\Lambda|}.
\]
For fixed $\beta$, the right side tends to zero.
\end{proof}

\subsection{A magnetic field as random extra pins}

The mixture in the next lemma is useful because its conditional laws
are precisely the deterministic pin laws already studied. Its mixing
measure need not be a product measure.

\begin{lemma}[Extra-pin mixture]
\label{lem:pin-mixture}
For $h\ge0$, the colored-picture law defining $Z^+_\Lambda(h)$ is
a mixture of the zero-field pinned laws with pin sets $E=D_0\cup F$,
where $F$ consists of time-zero slots. Under the mixing law,
\begin{equation}
 \Pr(i\in F)\le b(h):=1-e^{-\beta h}
 \qquad\text{for every time-zero slot }i.
 \label{eq:pin-marginal}
\end{equation}
Moreover,
\begin{equation}
 S-(p_N^+)'(h)
   =\frac1V\sum_{i\in V_\Lambda}\mathbb E_F[d_{D_0\cup F}(i)].
 \label{eq:field-deficit}
\end{equation}
\end{lemma}

\begin{proof}
For a colored picture let $D_{\rm down}$ count its down slots at time
zero. Then $M=SV-D_{\rm down}$. Start with independent selection of
all time-zero slots with probability $b(h)$. The probability that
only up slots are selected is
$(1-b(h))^{D_{\rm down}}=e^{-\beta hD_{\rm down}}$.
Consequently the joint weight obtained by multiplying the zero-field
colored-picture weight by the independent selection weight and by
$\ind{F\text{ contains only up slots}}$ has exactly the desired
field-tilted colored-picture marginal, after normalization.
Conditioning instead on $F$ leaves unit weight for each coloring
satisfying the pins $D_0\cup F$. Summing these colorings gives the
picture factor $2^{n_{D_0\cup F}}$, as required.

Conditional on the picture and its colors in this joint law, down
slots are never selected, while up slots are selected independently
with probability $b(h)$. This proves \eqref{eq:pin-marginal}.
Finally, differentiation of the finite-volume field tilt gives
$(p_N^+)'(h)=V^{-1}\mathbb E M$; conditioning on $F$ gives
\eqref{eq:field-deficit}.
\end{proof}

\subsection{The order of limits}

Recall the ideal heat-kernel constant and define its finite sums by
\begin{equation}
 \kappa=\frac1{(4\pi)^{3/2}\sqrt{\det D_{S,J}}},
 \qquad c_j=\kappa\sum_{n=1}^j n^{-3/2}.
 \label{eq:partial-bloch}
\end{equation}
Fix an integer $K\ge1$ and use the radius
$R=R_K=\lceil a(K+1)\beta\rceil$ from \eqref{eq:scales}.
The upper down-probability bound of Theorem~\ref{thm:return-bounds}
and \eqref{eq:field-deficit} imply
\begin{equation}
 S-(p_N^+)'(h)
 \le\beta^{-3/2}
       \bigl(c_{2K-1}+CK^{-1/2}+\epsilon_K(\beta)\bigr),
 \label{eq:pressure-upper}
\end{equation}
where $\epsilon_K(\beta)\ge0$ tends to zero as $\beta\to\infty$ for
fixed $K$, uniformly in $N$ and $h\ge0$. The leading constant $C$
is independent of $K$. There are $\ell V$ slots in the sum, so the
factor $1/\ell$ in each ideal slot-return probability has canceled.

For a lower bound, let $\rho_N$ be the fraction of sites $x\in\Lambda_N$
for which $\mathfrak B_R(\bar x)\subseteq\Lambda_N^\circ$. Since the
coset representatives and $A$ are fixed, all sites in this ball lie
within a bounded physical distance of $x$ when $R$ is fixed. Hence
$\rho_N\to1$ as $N\to\infty$, with $\beta,K$ fixed.
There are at most
\[
 v_R=\ell|\mathfrak C|(2R+1)^3
\]
time-zero slots above the ball. The union bound and
\eqref{eq:pin-marginal} show that it contains no extra pin with
probability at least $[1-v_Rb(h)]_+$, where $[u]_+=\max\{u,0\}$.
On this event Theorem~\ref{thm:return-bounds} supplies its lower bound.
Enlarging $\epsilon_K$ to cover both errors gives
\begin{equation}
 S-(p_N^+)'(h)
 \ge\rho_N[1-v_Rb(h)]_+\,
       \beta^{-3/2}\bigl(c_K-\epsilon_K(\beta)\bigr).
 \label{eq:pressure-lower}
\end{equation}
Here and below $\beta$ is sufficiently large, depending on $K$, that
$c_K-\epsilon_K(\beta)\ge0$.

\begin{proof}[Proof of Theorem~\ref{thm:main}]
Integrate \eqref{eq:pressure-upper} and \eqref{eq:pressure-lower}
from $0$ to $h>0$, and divide by $h$. This bounds
\[
 S-\frac{p_N^+(h)-p_N^+(0)}{h}.
\]
First send $N\to\infty$, using Proposition~\ref{prop:pressure-equivalence}
at both endpoints and $\rho_N\to1$. Then send $h\downarrow0$.
For these fixed $\beta,K$, continuity of $b$ at zero gives
\[
 \lim_{h\downarrow0}\frac1h\int_0^h[1-v_Rb(u)]_+\,du=1.
\]
By the definition of the right derivative,
\begin{equation}
 c_K-\epsilon_K(\beta)
 \le\beta^{3/2}\bigl(S-m_{S,J}(\beta)\bigr)
 \le c_{2K-1}+CK^{-1/2}+\epsilon_K(\beta).
 \label{eq:final-squeeze}
\end{equation}
Thus the thermodynamic limit and the zero-field derivative have been
handled through pressure increments. Now let $\beta\to\infty$ with
$K$ fixed and then $K\to\infty$. Both sides of
\eqref{eq:final-squeeze} converge to $\kappa\zeta(3/2)$, proving
the stated law.
\end{proof}
